% =====================================================================
% EN ESPERA. Bloques retirados de tablas.tex el 5-oct-2026 para
% reintroducirlos uno a uno, tras discutir si cada uno se sostiene.
% Este archivo no se compila. Para reintroducir un bloque: devolver su
% \input/figure a tablas.tex (antes de t12_heterogeneidad), su tabla de
% apendice y su prosa al Overview o a la Empirical Strategy.
% =====================================================================

% ===== ROBUSTEZ =====

% ---- 1. Placebo con outcomes fijados al nacer
% Prosa retirada:
%   Overview: The estimates are flat in a placebo on five outcomes fixed
%   at birth (Table~\ref{tab:placebo})
\begin{tabular}{l rl rl rl c}
\toprule
 & \multicolumn{2}{c}{1 a\~no$\times$EQ} & \multicolumn{2}{c}{2 a\~nos$\times$EQ} & \multicolumn{2}{c}{3--4 a\~nos$\times$EQ} & n \\
\midrule
Peso al nacer (z) & -0.013 & (0.067) & -0.082 & (0.080) & 0.017 & (0.075) & 9,153 \\
Talla al nacer (z) & -0.029 & (0.074) & -0.074 & (0.074) & 0.014 & (0.060) & 9,104 \\
Semanas de gestaci\'on (z) & -0.023 & (0.079) & -0.166** & (0.069) & -0.001 & (0.077) & 9,715 \\
Parto prematuro (0/1) & 0.017 & (0.021) & 0.041* & (0.024) & 0.003 & (0.021) & 9,988 \\
Apgar 5 min.\ (z) & -0.083 & (0.087) & -0.043 & (0.089) & -0.013 & (0.112) & 4,143 \\
\bottomrule
\end{tabular}

% ---- begin paper/tables/a5_nacimiento.tex
\begin{table}[htbp]\centering
\caption{Placebo Outcomes Fixed at Birth: Items and Scales}
\label{tab:a_birth}
\footnotesize\renewcommand{\arraystretch}{1.0}
\begin{threeparttable}
\begin{tabular}{@{}p{68pt}p{212pt}p{150pt}@{}}
\toprule
 & Question, verbatim & Response scale \\
 & (1) & (2) \\
\midrule
\texttt{g24} & ?`Cuál fue el peso del(de la) niño(a) al nacer? & Open numeric answer (kilograms); special codes: 88 = dato impugnado; 99 = no sabe no responde \\
\addlinespace[1pt]
\texttt{g23} & ?`Cuál fue la talla, altura del niño(a) al nacer? & Open numeric answer (centimeters); special codes: 88 = dato impugnado; 99 = no sabe no responde \\
\addlinespace[1pt]
\texttt{g20} & ?`Cuántas semanas de embarazo tenía Ud. cuando nació el(la) niño(a)? & Open numeric answer (weeks); special codes: 88 = dato impugnado; 99 = no sabe no responde \\
\addlinespace[1pt]
\texttt{g17\_6} & Parto prematuro (antes de las 37 semanas). Al momento del nacimiento, ?`ex\,\dots & 1 = sí; 2 = no; 9 = no sabe no responde \\
\addlinespace[1pt]
\texttt{g25b1} & b1)Medición 2 (5 Minutos). ?`Cuál fue su puntuación en el Test de Apgar? & Open numeric answer (score 0--10); special codes: 99 = no sabe no responde \\
\bottomrule
\end{tabular}
\begin{wpnotes}
\textit{Notes:} Question wording and response labels are reproduced verbatim in the original Spanish of the questionnaire, from the variable and value labels of the public microdata; wording the microdata truncates at 80 characters is marked with an ellipsis. Cleaning: birth weight kept between 0.5 and 6.5 kilograms, length between 30 and 65 centimeters, gestation between 24 and 45 weeks, Apgar between 0 and 10; the survey codes 99 for a forgotten answer, treated as missing. All five outcomes are fixed at birth, one to four years before the earthquake for every cohort of the analysis sample.
\end{wpnotes}
\end{threeparttable}
\end{table}
% ---- end paper/tables/a5_nacimiento.tex


% ---- 2. Atricion 2010-2024
% ---- begin paper/tables/t9_atricion.tex
\begin{table}[htbp]\centering
\caption{Attrition from 2010 to 2024 Is Not Differential by Exposure Age and Zone}
\label{tab:attrition}
\begin{threeparttable}
\begin{tabular}{@{}p{250pt}>{\centering\arraybackslash}p{130pt}@{}}
\toprule
 & Pr(reinterviewed in 2024) \\
 & (1) \\
\midrule
Exposed at age 1 $\times$ affected & 0.038 \\
 & (0.029) \\
\addlinespace
Exposed at age 2 $\times$ affected & 0.028 \\
 & (0.029) \\
\addlinespace
Exposed at ages 3--4 $\times$ affected & 0.038 \\
 & (0.029) \\
\midrule
Exposure-age and region fixed effects & Yes \\
Mother's education and age in 2010 & Yes \\
Observations & 14,855 \\
Mean retention into 2024 & 67.3 percent \\
\bottomrule
\end{tabular}
\begin{wpnotes}
\textit{Notes:} Linear probability model over the children of the 2010 baseline with a valid birth date. The outcome is an indicator for being reinterviewed in the 2024 wave. The small and insignificant coefficients show that the loss of the sample (32.7 percent) does not concentrate in any treatment-by-age cell; the fitted probabilities feed the IPW row of Table~\ref{tab:robust}. The affected zone enters through the region of residence reported in 2010. Robust standard errors in parentheses. \starnote
\end{wpnotes}
\end{threeparttable}
\end{table}
% ---- end paper/tables/t9_atricion.tex


% ---- 3. Robustez de especificacion e inferencia exigente
% Prosa retirada:
%   Overview: and stable across every specification we tried
%   (Table~\ref{tab:robust}). Permutation of the affected zone across the
%   116 municipalities gives $p=0.045$
%   Strategy: and Table~\ref{tab:robust} adds permutation and wild-cluster
%   inference for the preferred specification.
%   Nota de T3 (scripts/15, t3_main): Demanding inference for column 3 is
%   reported in Table~\ref{tab:robust}.
% ---- begin paper/tables/t7_robustez.tex
\begin{table}[htbp]\centering
\caption{Robustness and Demanding Inference for the Main Result}
\label{tab:robust}
\footnotesize\renewcommand{\arraystretch}{0.95}
\begin{threeparttable}
\begin{tabular}{@{}p{215pt}*{1}{>{\centering\arraybackslash}p{120pt}}>{\centering\arraybackslash}p{62pt}@{}}
\toprule
 & Dependent variable: PHQ-4 score (z) & \\
\cmidrule(lr){2-2}
 & Age 2 $\times$ affected & Observations \\
 & (1) & (2) \\
\midrule
Preferred specification (column 3) & $-$0.147\sym{**} & 9,996 \\
 & (0.069) & \\
\addlinespace
Fixed effects only & $-$0.144\sym{**} & 9,996 \\
 & (0.072) & \\
\addlinespace
Excluding the Metropolitan Region & $-$0.120 & 6,430 \\
 & (0.076) & \\
\addlinespace
Unweighted & $-$0.129\sym{*} & 9,996 \\
 & (0.071) & \\
\addlinespace
PHQ-4 in points (0--12), not standardized & $-$0.391\sym{**} & 9,996 \\
 & (0.183) & \\
\addlinespace
Balanced sample with a 2017 CBCL & $-$0.114 & 7,188 \\
 & (0.082) & \\
\addlinespace
Controlling for the birth endowment & $-$0.144\sym{**} & 9,996 \\
 & (0.069) & \\
\addlinespace
Reweighted for attrition (IPW) & $-$0.149\sym{**} & 9,996 \\
 & (0.069) & \\
\addlinespace
Continuous dose: PGA (z) instead of the zone & $-$0.068\sym{**} & 9,996 \\
 & (0.027) & \\
\addlinespace
Dropping one affected region at a time (range) & [$-$0.172, $-$0.120] & \\
Municipality and birth-cohort fixed effects (all rows) & Yes & \\
\midrule
\panel{3}{Demanding inference for the preferred specification ($p$-values)}
\hspace{1em}Permutation of the affected zone & 0.045 & \\
\hspace{1em}Wild cluster bootstrap by region (15 regions) & 0.132 & \\
\bottomrule
\end{tabular}
\begin{wpnotes}
\textit{Notes:} Each row re-estimates the preferred specification of Table~\ref{tab:main}, Panel A, changing one thing. Column 1 reports the coefficient on exposure at age two $\times$ affected zone. The IPW row reweights by the inverse of the estimated probability of remaining in the 2024 wave (Table~\ref{tab:attrition}). The permutation test reassigns the affected-zone status across the 116 municipalities of selection, keeping the number of treated municipalities, 2,000 times. The wild cluster bootstrap by region uses Webb weights and clusters at the level of the regional treatment. \clusternote{} \starnote
\end{wpnotes}
\end{threeparttable}
\end{table}
% ---- end paper/tables/t7_robustez.tex


% ---- 4. Romano-Wolf (familia de cuatro outcomes)
% Prosa retirada:
%   Overview: and adjusting the four mental-health outcomes as a family
%   leaves Romano--Wolf $p$-values between 0.10 and 0.12, so we read the
%   full-sample evidence as suggestive and the girls' result as the firm
%   finding.
\begin{tabular}{l cc cc}
\toprule
 & \multicolumn{2}{c}{2 a\~nos$\times$EQ} & \multicolumn{2}{c}{Contraste 3--4$-$2} \\
 & $p$ anal\'itico & $p$ RW & $p$ anal\'itico & $p$ RW \\
\midrule
PHQ-4 (z) & 0.033 & 0.120 & 0.088 & 0.235 \\
PHQ-2 positivo & 0.247 & 0.299 & 0.244 & 0.468 \\
GAD-2 positivo & 0.022 & 0.102 & 0.032 & 0.111 \\
CBCL 2024 (z) & 0.147 & 0.299 & 0.583 & 0.615 \\
\bottomrule
\end{tabular}


% ---- 5. Dosis sismica continua (PGA) y su mapa
% Prosa retirada:
%   Strategy (tras Astroza et al. 2010): and the dose tables replace it with
%   the standardized peak ground acceleration, the strength with which the
%   ground actually shook in the municipality;
% ---- begin paper/tables/t5_dosis.tex
\begin{table}[htbp]\centering
\caption{The Gradient Under Continuous Seismic Doses}
\label{tab:dose}
\begin{threeparttable}
\begin{tabular}{@{}p{170pt}*{4}{>{\centering\arraybackslash}p{67pt}}@{}}
\toprule
 & \multicolumn{4}{c}{Seismic dose of the municipality of selection} \\
\cmidrule(lr){2-5}
 & Affected zone (0/1) & PGA (z) & MMI (z) & $-$log dist.\ to epicenter (z) \\
 & (1) & (2) & (3) & (4) \\
\midrule
\panel{5}{Panel A. PHQ-4 score: sum of four depression--anxiety items (z)}
Exposed at 12--23 & $-$0.045 & $-$0.035 & $-$0.028 & $-$0.029 \\
 & (0.078) & (0.033) & (0.032) & (0.029) \\
\addlinespace[3pt]
Exposed at 24--35 & $-$0.147\sym{**} & $-$0.068\sym{**} & $-$0.062\sym{***} & $-$0.062\sym{**} \\
 & (0.069) & (0.027) & (0.023) & (0.025) \\
\addlinespace[3pt]
Exposed at 36--59 & $-$0.044 & $-$0.037 & $-$0.039 & $-$0.025 \\
 & (0.060) & (0.027) & (0.024) & (0.023) \\
\addlinespace
\panel{5}{Panel B. Positive GAD-2 screen: two anxiety items, subscore of 3 or more}
Exposed at 12--23 & $-$0.012 & $-$0.015 & $-$0.013 & $-$0.022\sym{*} \\
 & (0.033) & (0.014) & (0.013) & (0.013) \\
\addlinespace[3pt]
Exposed at 24--35 & $-$0.073\sym{**} & $-$0.039\sym{***} & $-$0.029\sym{**} & $-$0.036\sym{***} \\
 & (0.032) & (0.015) & (0.013) & (0.013) \\
\addlinespace[3pt]
Exposed at 36--59 & $-$0.016 & $-$0.016 & $-$0.010 & $-$0.014 \\
 & (0.028) & (0.013) & (0.010) & (0.011) \\
\midrule
Municipality fixed effects & Yes & Yes & Yes & Yes \\
Birth-cohort fixed effects & Yes & Yes & Yes & Yes \\
Pre-earthquake controls & Yes & Yes & Yes & Yes \\
2012 controls & Yes & Yes & Yes & Yes \\
Adolescents & 9,996 & 9,996 & 9,996 & 9,996 \\
Selection municipalities (clusters) & 116 & 116 & 116 & 116 \\
\bottomrule
\end{tabular}
\begin{wpnotes}
\textit{Notes:} Estimates of equation~(1) under four measures of the seismic dose, all with the controls of column 3 of Table~\ref{tab:main}. Column 1 repeats the official binary affected zone. Column 2 uses the peak ground acceleration: the strongest shaking of the ground during the earthquake, expressed as a share of the acceleration of gravity, measured at the centroid of the municipality of selection (from the ShakeMap of the U.S. Geological Survey) and standardized; column 3 an intensity scale from the same source; column 4 minus the log of the distance from the centroid to the epicenter. The continuous doses vary across the 116 municipalities, within affected regions, which addresses inference concerns with a regional treatment. \bindef{} \clusternote{} \starnote
\end{wpnotes}
\end{threeparttable}
\end{table}
% ---- end paper/tables/t5_dosis.tex


\begin{figure}[htbp]\centering
  \includegraphics[width=0.9\textwidth]{f2_dosis_pga.pdf}
  \caption{The Gradient Under the Continuous Seismic Dose}
  \label{fig:pga}
  \begin{fignotes}[0.9\textwidth]
\textit{Notes:} Each marker is the coefficient on the exposure-age group
interacted with the standardized peak ground acceleration, how strongly the
ground shook, at the centroid of the municipality of selection, with
its 95 percent confidence interval; specification of column 3 of
Table~\ref{tab:main}. The dose varies across the 116 municipalities, within affected
regions. \clusternote
  \end{fignotes}
\end{figure}

\begin{figure}[htbp]\centering
  \includegraphics[height=0.86\textheight]{f3_mapa_dosis.pdf}
  \caption{The Seismic Dose Across the Municipalities of Selection}
  \label{fig:mapdose}
  \begin{fignotes}[0.86\textwidth]
\textit{Notes:} Each colored municipality is one of the 116 municipalities of
selection of the ELPI sample, colored by the peak ground acceleration, that
is, how strongly the ground shook during the earthquake, expressed as a share
of the acceleration of gravity (from the ShakeMap of the U.S. Geological
Survey); the star marks the epicenter and light-gray municipalities are not
in the sample. The map shows the two sources of identifying variation of
Table~\ref{tab:dose}, the affected zone and the continuous dose within it.
  \end{fignotes}
\end{figure}

% ---- 6. Mapa descriptivo de ansiedad 2024
\begin{figure}[htbp]\centering
  \includegraphics[height=0.86\textheight]{f4_mapa_sintomas.pdf}
  \caption{The Geography of Adolescent Anxiety in 2024}
  \label{fig:mapanx}
  \begin{fignotes}[0.86\textwidth]
\textit{Notes:} Each colored municipality is one of the 116 municipalities of
selection, colored by the share of its adolescents with a positive GAD-2
anxiety screen in the 2024 wave; municipality shares range from 11.9 to 43.8
percent, and light-gray municipalities are not in the sample. The map is
descriptive and unadjusted; the regressions of the paper compare cohorts
within municipality, never municipalities against each other.
  \end{fignotes}
\end{figure}

% ===== MECANISMOS =====

% ---- 7. Otros resultados adolescentes I: bienestar y victimizacion
% ---- begin paper/tables/t10a_bienestar.tex
\begin{table}[htbp]\centering
\caption{Other Adolescent Outcomes I: Well-Being and Victimization}
\label{tab:mech1}
\begin{threeparttable}
\begin{tabular}{@{}p{148pt}*{5}{>{\centering\arraybackslash}p{62pt}}@{}}
\toprule
 & Resilience (z) & Life satisf. (z) & Poor health (z) & Bullying (z) & Cyber-vict. (0/1) \\
 & (1) & (2) & (3) & (4) & (5) \\
\midrule
Exposed at age 1 & 0.061 & 0.086\sym{*} & $-$0.086 & 0.028 & $-$0.015 \\
 & (0.077) & (0.050) & (0.066) & (0.060) & (0.029) \\
\addlinespace[3pt]
Exposed at age 2 & 0.141\sym{*} & 0.098\sym{*} & $-$0.072 & $-$0.047 & $-$0.007 \\
 & (0.079) & (0.056) & (0.077) & (0.071) & (0.023) \\
\addlinespace[3pt]
Exposed at ages 3--4 & 0.098 & 0.116\sym{*} & $-$0.054 & 0.037 & $-$0.005 \\
 & (0.068) & (0.062) & (0.070) & (0.063) & (0.025) \\
\midrule
Municipality fixed effects & Yes & Yes & Yes & Yes & Yes \\
Birth-cohort fixed effects & Yes & Yes & Yes & Yes & Yes \\
Pre-earthquake controls & Yes & Yes & Yes & Yes & Yes \\
2012 controls & Yes & Yes & Yes & Yes & Yes \\
Observations & 9,996 & 9,996 & 9,996 & 9,176 & 9,985 \\
\bottomrule
\end{tabular}
\begin{wpnotes}
\textit{Notes:} Each column reports one estimate of equation~(1) with the specification of column 3 of Table~\ref{tab:main} on another outcome of the 2024 wave. Resilience is the mean of the six Brief Resilience Scale items; life satisfaction a single 1--7 item; poor health a single reversed item; bullying the mean of eight frequency items, one reverse-coded; cyber-victimization one if any of six items is answered yes. Appendix Table~\ref{tab:a_other} lists every item and scale verbatim. \bindef{} \clusternote{} \starnote
\end{wpnotes}
\end{threeparttable}
\end{table}
% ---- end paper/tables/t10a_bienestar.tex

% (a2_wellbeing y a3_conducta ya estan en el apendice principal)

% ---- 8. Otros resultados adolescentes II: violencia de pareja y sustancias
% ---- begin paper/tables/t10b_riesgo.tex
\begin{table}[htbp]\centering
\caption{Other Adolescent Outcomes II: Dating Violence and Substance Use}
\label{tab:mech2}
\begin{threeparttable}
\begin{tabular}{@{}p{148pt}*{4}{>{\centering\arraybackslash}p{62pt}}@{}}
\toprule
 & Dating viol. (0/1) & Tobacco (0/1) & Alcohol (0/1) & Cannabis (0/1) \\
 & (1) & (2) & (3) & (4) \\
\midrule
Exposed at age 1 & 0.025 & 0.012 & $-$0.038 & 0.018 \\
 & (0.019) & (0.016) & (0.032) & (0.017) \\
\addlinespace[3pt]
Exposed at age 2 & 0.005 & $-$0.002 & $-$0.037 & 0.014 \\
 & (0.020) & (0.015) & (0.040) & (0.016) \\
\addlinespace[3pt]
Exposed at ages 3--4 & 0.028 & $-$0.022 & $-$0.078\sym{**} & 0.013 \\
 & (0.021) & (0.024) & (0.031) & (0.019) \\
\midrule
Municipality fixed effects & Yes & Yes & Yes & Yes \\
Birth-cohort fixed effects & Yes & Yes & Yes & Yes \\
Pre-earthquake controls & Yes & Yes & Yes & Yes \\
2012 controls & Yes & Yes & Yes & Yes \\
Observations & 7,101 & 9,995 & 9,995 & 9,995 \\
\bottomrule
\end{tabular}
\begin{wpnotes}
\textit{Notes:} Each column reports one estimate of equation~(1) with the specification of column 3 of Table~\ref{tab:main} on another outcome of the 2024 wave. Dating violence is one if any of three items is answered yes, among adolescents who ever dated; tobacco, alcohol and cannabis refer to the last 12 months, with alcohol one if any of four beverage items is answered yes. Appendix Table~\ref{tab:a_other} lists every item and scale verbatim. \bindef{} \clusternote{} \starnote
\end{wpnotes}
\end{threeparttable}
\end{table}
% ---- end paper/tables/t10b_riesgo.tex


% ---- 9. Salud mental previa de la madre
\begin{tabular}{l rl rl rl rl}
\toprule
 & \multicolumn{4}{c}{PHQ-4 (z)} & \multicolumn{4}{c}{GAD-2 positivo} \\
\cmidrule(lr){2-5} \cmidrule(lr){6-9}
 & \multicolumn{2}{c}{Madre vulnerable} & \multicolumn{2}{c}{Madre no vulnerable} & \multicolumn{2}{c}{Madre vulnerable} & \multicolumn{2}{c}{Madre no vulnerable} \\
\midrule
Expuesto con 1 a\~no (12--23 m.) & -0.036 & (0.174) & -0.021 & (0.080) & -0.022 & (0.066) & -0.001 & (0.036) \\
Expuesto con 2 a\~nos (24--35 m.) & -0.101 & (0.129) & -0.124 & (0.089) & -0.061 & (0.063) & -0.062 & (0.040) \\
Expuesto con 3--4 a\~nos (36--59 m.) & -0.003 & (0.127) & -0.017 & (0.071) & -0.052 & (0.063) & 0.010 & (0.031) \\
\addlinespace
Contraste 3--4 $-$ 2 & 0.097 & (0.169) & 0.107 & (0.070) & 0.008 & (0.065) & 0.072** & (0.031) \\
\midrule
Adolescentes & \multicolumn{2}{c}{2,023} & \multicolumn{2}{c}{7,973} & \multicolumn{2}{c}{2,023} & \multicolumn{2}{c}{7,973} \\
\bottomrule
\end{tabular}

% (a4_madre ya esta en el apendice principal; si vuelve el bloque,
% recuperar su Panel A de items 2010 de embarazo y posparto)

